\documentclass[11pt]{article}

\usepackage[a4paper,margin=1in]{geometry}
\usepackage{amsmath,amssymb,amsthm,mathtools,bm}
\usepackage[T1]{fontenc}
\usepackage{lmodern}
\usepackage{microtype}
\usepackage{enumitem}
\usepackage{hyperref}

\allowdisplaybreaks
\numberwithin{equation}{section}

\newtheorem{theorem}{Theorem}[section]
\newtheorem{proposition}[theorem]{Proposition}
\newtheorem{lemma}[theorem]{Lemma}
\newtheorem{corollary}[theorem]{Corollary}
\newtheorem{definition}[theorem]{Definition}
\newtheorem{remark}[theorem]{Remark}
\newtheorem{gap}[theorem]{Gap}

\title{Current State of the Connectedness Threshold Proof\\
for the Open-Boundary Hatano--Nelson Matrix}
\author{}
\date{}

\begin{document}
\maketitle

\begin{abstract}
This note records the current partial proof of the connectedness threshold problem
for
\[
\mathbf A_n(a)=\operatorname{tridiag}(a,0,1),\qquad 0<a<1,\qquad n\ge2.
\]
The uploaded manuscript proves that connectedness of the pseudospectrum is
equivalent to crossing a one-dimensional real-axis barrier. Thus the threshold
problem is reduced to proving
\[
\eta_{n+1}(a)\le \eta_n(a),\qquad n\ge2.
\]
The present note records the additional analytic and Weyl-contact reductions
developed so far. In particular, it proves the endpoint and large-\(n\) exclusions
used to reduce the problem to a finite middle block, derives the reversal-pencil
active-contact identities, and records the elimination of the exact-rigid
even-to-odd Weyl branch. The note ends with the currently unfinished gaps.
\end{abstract}

\tableofcontents

\section{Basic notation and the barrier reduction}

Fix
\[
0<a<1,\qquad n\ge2,
\]
and define
\[
\mathbf A_n(a)=
\begin{bmatrix}
0&1\\
a&0&1\\
&\ddots&\ddots&\ddots\\
&&a&0&1\\
&&&a&0
\end{bmatrix}.
\]
For \(z=x+iy\), write
\[
s_n(x,y)
=
s_{\min}\!\bigl((x+iy)\mathbf I_n-\mathbf A_n(a)\bigr),
\qquad
s_n(x):=s_n(x,0).
\]
The \(\varepsilon\)-pseudospectrum is
\[
\Sigma_\varepsilon^{(n)}
=
\{z\in\mathbb C:\ s_n(\Re z,\Im z)<\varepsilon\}.
\]
The real slice is
\[
I_\varepsilon^{(n)}
=
\Sigma_\varepsilon^{(n)}\cap\mathbb R
=
\{x\in\mathbb R:\ s_n(x)<\varepsilon\}.
\]

The eigenvalues of \(\mathbf A_n(a)\) are
\[
\mu_k^{(n)}
=
2\sqrt a\cos\frac{k\pi}{n+1},
\qquad k=1,\dots,n,
\]
and
\[
\mu_1^{(n)}>\mu_2^{(n)}>\cdots>\mu_n^{(n)}.
\]
Define the \(k\)-th real-axis barrier by
\[
\beta_{n,k}(a)
=
\max_{x\in[\mu_{k+1}^{(n)},\mu_k^{(n)}]}s_n(x),
\qquad k=1,\dots,n-1,
\]
and define the global real-axis barrier
\[
\eta_n(a):=\max_{1\le k\le n-1}\beta_{n,k}(a).
\]

\begin{theorem}[Barrier criterion, imported from the manuscript]
\label{thm:barrier}
For every \(n\ge2\) and every \(\varepsilon>0\),
\[
\Sigma_\varepsilon^{(n)}\text{ is connected}
\quad\Longleftrightarrow\quad
I_\varepsilon^{(n)}\text{ is connected}
\quad\Longleftrightarrow\quad
\varepsilon>\eta_n(a).
\]
\end{theorem}

\begin{corollary}[Threshold reduction]
\label{cor:threshold-reduction}
The connectedness threshold theorem follows from
\[
\eta_{n+1}(a)\le\eta_n(a),
\qquad n\ge2,\quad 0<a<1.
\]
Indeed, if this monotonicity holds, then
\[
\Sigma_\varepsilon^{(n)}\text{ connected}
\quad\Longleftrightarrow\quad
\eta_n(a)<\varepsilon
\]
is an upper-tail condition in \(n\).
\end{corollary}

Thus the remaining connectedness threshold problem is the monotonicity problem
\[
\boxed{
\eta_{n+1}(a)\le\eta_n(a).
}
\]

\section{Analytic exclusions outside a finite middle block}

This section records the analytic exclusions currently available. They do not yet
prove the theorem for all \(n,a\), but they imply that any counterexample must
belong to a finite middle block, provided the asymptotic estimates stated below are
accepted as established.

\subsection{\texorpdfstring{Uniform small-\(a\) exclusion}{Uniform small-a exclusion}}

Let
\[
r=\sqrt a.
\]
Write
\[
\mathbf A_n(a)=S_n+r^2S_n^T,
\]
where \(S_n\) is the upper shift. For
\[
x=2r\cos\theta,
\]
set
\[
B_n(\theta)=2r\cos\theta\,\mathbf I_n-\mathbf A_n(a).
\]
The real spectral hull corresponds to
\[
\theta\in\left[\frac{\pi}{n+1},\frac{n\pi}{n+1}\right].
\]

\begin{lemma}[Uniform small-\(a\) upper bound]
\label{lem:small-upper}
If \(0<r\le1/10\), then
\[
\eta_{n+1}(a)\le r^{n+1}(n+2).
\]
\end{lemma}

\begin{proof}
For \(\theta\) not on the eigenvalue grid, the tridiagonal Toeplitz inverse formula gives
\[
(B_n(\theta)^{-1})_{ij}
=
\frac{r^{i-j-1}}{\sin\theta\,\sin((n+1)\theta)}
\begin{cases}
\sin(i\theta)\sin((n+1-j)\theta),&i\le j,\\
\sin(j\theta)\sin((n+1-i)\theta),&i\ge j.
\end{cases}
\]
Let \(\Pi_n\) be the reversal matrix. The \((1,1)\)-entry of
\(B_n(\theta)^{-1}\Pi_n\) is
\[
\frac{r^{-n}\sin\theta}{\sin((n+1)\theta)}.
\]
Thus
\[
\|B_n(\theta)^{-1}\|_2
\ge
r^{-n}\frac{\sin\theta}{|\sin((n+1)\theta)|}.
\]
Therefore
\[
s_n(2r\cos\theta)
\le
r^n\frac{|\sin((n+1)\theta)|}{\sin\theta}
=
r^n|U_n(\cos\theta)|,
\]
where \(U_n\) is the Chebyshev polynomial of the second kind. Since
\[
|U_n(\cos\theta)|\le n+1,
\]
we get
\[
\eta_n(a)\le r^n(n+1),
\]
and replacing \(n\) by \(n+1\) gives the claim.
\end{proof}

\begin{lemma}[Uniform small-\(a\) lower bound]
\label{lem:small-lower}
For \(0<r\le1/10\) and every \(n\ge2\),
\[
\eta_n(a)\ge \frac{r^n(n+1)}{C_*},
\qquad
C_*<5.3.
\]
\end{lemma}

\begin{proof}
Choose
\[
\theta_n=\frac{3\pi}{2(n+1)}.
\]
Then
\[
\theta_n\in\left[\frac{\pi}{n+1},\frac{n\pi}{n+1}\right],
\qquad
|\sin((n+1)\theta_n)|=1.
\]
Using the inverse formula above and multiplying on the right by \(\Pi_n\), one writes
\[
B_n(\theta_n)^{-1}\Pi_n=H_n^{(1)}+H_n^{(2)},
\]
where \(H_n^{(1)}\) is the rank-one dominant triangular part and
\(H_n^{(2)}\) is the reflected triangular remainder. The dominant part obeys
\[
\|H_n^{(1)}\|_2
\le
1.05\frac{3\pi}{2}\frac{r^{-n}}{n+1},
\]
because
\[
\sum_{j=1}^\infty j^2r^{2(j-1)}
=
\frac{1+r^2}{(1-r^2)^3}
<1.05
\qquad (r\le1/10).
\]
The reflected remainder satisfies
\[
\|H_n^{(2)}\|_2
\le
\frac13\frac{r^{-n}}{n+1}.
\]
Hence
\[
\|B_n(\theta_n)^{-1}\|_2
\le
C_*\frac{r^{-n}}{n+1},
\qquad
C_*=
1.05\frac{3\pi}{2}+\frac13<5.3.
\]
Therefore
\[
s_n(2r\cos\theta_n)
=
\|B_n(\theta_n)^{-1}\|_2^{-1}
\ge
\frac{r^n(n+1)}{C_*}.
\]
Since \(\eta_n(a)\) is the maximum over the real hull, the result follows.
\end{proof}

\begin{theorem}[Uniform small-\(a\) monotonicity]
\label{thm:small-a}
For
\[
0<a\le10^{-2},
\]
one has
\[
\eta_{n+1}(a)<\eta_n(a),
\qquad n\ge2.
\]
\end{theorem}

\begin{proof}
By Lemmas~\ref{lem:small-upper} and \ref{lem:small-lower},
\[
\frac{\eta_{n+1}(a)}{\eta_n(a)}
\le
rC_*\frac{n+2}{n+1}.
\]
For \(n\ge2\),
\[
\frac{n+2}{n+1}\le\frac43.
\]
If \(r\le1/10\), then
\[
rC_*\frac{n+2}{n+1}
\le
\frac1{10}\cdot5.3\cdot\frac43<1.
\]
Thus \(\eta_{n+1}(a)<\eta_n(a)\).
\end{proof}

\subsection{Perturbative near-Hermitian exclusion}

At \(a=1\), the matrix
\[
\mathbf A_n(1)=\operatorname{tridiag}(1,0,1)
\]
is Hermitian. Hence
\[
s_n(x)=\operatorname{dist}(x,\sigma(\mathbf A_n(1))).
\]
The eigenvalues are
\[
\lambda_k^{(n)}=2\cos\frac{k\pi}{n+1}.
\]
Thus \(\eta_n(1)\) is half of the largest spectral gap. Explicitly,
\[
\eta_n(1)=
\begin{cases}
2\sin\frac{\pi}{2(n+1)},&n\text{ even},\\[1mm]
\sin\frac{\pi}{n+1},&n\text{ odd}.
\end{cases}
\]

\begin{lemma}[Hermitian gap]
\label{lem:hermitian-gap}
For \(n\ge2\),
\[
\eta_n(1)-\eta_{n+1}(1)
\ge
\frac{1}{(n+2)^2}.
\]
\end{lemma}

\begin{proof}
This follows directly from the preceding trigonometric formulas and the mean value
theorem applied to \(\sin t\) on \([0,\pi/2]\). The constant is deliberately conservative.
\end{proof}

\begin{theorem}[Inner near-Hermitian layer]
\label{thm:inner-layer}
If
\[
1-a\le\frac{1}{4(n+2)^2},
\]
then
\[
\eta_{n+1}(a)<\eta_n(a).
\]
\end{theorem}

\begin{proof}
Since
\[
\mathbf A_n(a)-\mathbf A_n(1)=(a-1)S_n^T,
\]
we have
\[
\|\mathbf A_n(a)-\mathbf A_n(1)\|_2=1-a.
\]
The least singular value is \(1\)-Lipschitz in the matrix, hence
\[
|\eta_n(a)-\eta_n(1)|\le1-a
\]
provided the Hermitian maximizer remains inside the smaller real spectral hull.
Under the displayed hypothesis this is true, because \(a>3/4\). Therefore
\[
\eta_n(a)-\eta_{n+1}(a)
\ge
\eta_n(1)-\eta_{n+1}(1)-2(1-a).
\]
Using Lemma~\ref{lem:hermitian-gap},
\[
\eta_n(a)-\eta_{n+1}(a)
\ge
\frac{1}{(n+2)^2}-2(1-a)
\ge
\frac{1}{2(n+2)^2}>0.
\]
\end{proof}

\subsection{\texorpdfstring{Large-\(n\) middle and near-Hermitian transition estimates}{Large-n middle and near-Hermitian transition estimates}}

The following estimates were derived from the exact Green kernel of
\[
B_n(\theta)=2r\cos\theta\,\mathbf I_n-\mathbf A_n(a),
\qquad r=\sqrt a.
\]
They are part of the current proof package, but their full polished write-up should
include all uniform remainder estimates.

\begin{theorem}[Compact-middle eventual monotonicity]
\label{thm:compact-middle}
For every compact interval
\[
0<a_-<a_+<1,
\]
there exists \(N=N(a_-,a_+)\) such that
\[
a\in[a_-,a_+],\quad n\ge N
\quad\Longrightarrow\quad
\eta_{n+1}(a)<\eta_n(a).
\]
\end{theorem}

\begin{proof}[Proof sketch]
Set \(r=\sqrt a\) and \(x=2r\cos\theta\). The exact inverse kernel gives, uniformly for
\(r\in[r_-,r_+]\subset(0,1)\),
\[
s_n(2r\cos\theta)
=
r^n|\sin((n+1)\theta)|
\frac{\sin\theta}{S_r(\theta)}
\left(1+O(n^3r_+^n)\right),
\]
where
\[
S_r(\theta)=\sum_{j=1}^{\infty}r^{2(j-1)}\sin^2(j\theta).
\]
Near the endpoint \(\theta=0\),
\[
\frac{\sin\theta}{S_r(\theta)}
=
\frac{(1-r^2)^3}{1+r^2}\frac1\theta+O(\theta).
\]
Therefore
\[
\eta_n(a)
=
r^n(n+1)
\frac{(1-r^2)^3}{1+r^2}
\kappa
\left(1+O\!\left(\frac1n\right)\right),
\]
uniformly on compact subintervals of \((0,1)\), where
\[
\kappa=\max_{u\ge\pi}\frac{|\sin u|}{u}.
\]
Thus
\[
\frac{\eta_{n+1}(a)}{\eta_n(a)}
=
r\frac{n+2}{n+1}\left(1+O\!\left(\frac1n\right)\right),
\]
uniformly for \(a\in[a_-,a_+]\). Since \(r\le\sqrt{a_+}<1\), the ratio is eventually
strictly less than \(1\).
\end{proof}

\begin{theorem}[Near-Hermitian mesoscopic transition estimate]
\label{thm:mesoscopic}
Let
\[
N=n+1,\qquad
a=e^{-2\Gamma/N},
\qquad
1\le\Gamma\le C\log N.
\]
Let
\[
(\mathcal V_\Gamma f)(s)
=
\int_s^1 e^{\Gamma(t-s)}f(t)\,dt
\]
and
\[
\sigma(\Gamma)=\|\mathcal V_\Gamma\|_2^{-1}.
\]
Then
\[
\eta_n(e^{-2\Gamma/N})
=
\frac{2\sigma(\Gamma)}{N}
+
\frac{E(\Gamma)}{N^2}
+
O_C\!\left(\frac{\sigma(\Gamma)(1+\Gamma)^2}{N^3}\right),
\]
uniformly for \(1\le\Gamma\le C\log N\), with
\[
|E(\Gamma)|+|E'(\Gamma)|
\le
K_C\sigma(\Gamma)(1+\Gamma).
\]
Consequently, for every \(C>0\), there exists \(N_C\) such that
\[
1\le -\frac{n+1}{2}\log a\le C\log(n+1),\quad n\ge N_C
\]
implies
\[
\eta_{n+1}(a)<\eta_n(a).
\]
\end{theorem}

\begin{proof}[Proof sketch]
The exact Green kernel, after staggering near the central gap
\[
\theta=\frac{\pi}{2}+O(N^{-1}),
\]
converges after scaling to the Volterra operator \(\mathcal V_\Gamma\). The
Euler--Maclaurin expansion of the kernel gives the displayed two-term expansion.
For fixed \(a\), set
\[
\Gamma_N(a)=-\frac{N}{2}\log a.
\]
Then
\[
\Gamma_{N+1}(a)=\Gamma_N(a)\left(1+\frac1N\right).
\]
The two-term expansion gives
\[
\eta_n(a)-\eta_{n+1}(a)
=
\frac{2(\sigma(\Gamma)-\Gamma\sigma'(\Gamma))}{N^2}
+
O_C\!\left(
\frac{\sigma(\Gamma)(1+\Gamma)^2}{N^3}
\right).
\]
The kernel of \(\mathcal V_\Gamma\) increases pointwise with \(\Gamma\), so
\(\|\mathcal V_\Gamma\|\) is strictly increasing and \(\sigma'(\Gamma)<0\). Hence
\[
\sigma(\Gamma)-\Gamma\sigma'(\Gamma)>\sigma(\Gamma)>0.
\]
Since \((1+\Gamma)^2/N\to0\) for \(\Gamma\le C\log N\), the difference is positive
for all large \(N\).
\end{proof}

\begin{corollary}[Finite middle block reduction]
\label{cor:finite-block}
Assuming Theorems~\ref{thm:compact-middle} and \ref{thm:mesoscopic} with their
uniform remainders, any counterexample to
\[
\eta_{n+1}(a)\le\eta_n(a)
\]
must lie in a finite set of sizes
\[
2\le n<N_*,
\]
and in a compact interval
\[
a\in[a_0,1-\delta]\subset(0,1).
\]
\end{corollary}

\section{The symmetric reversal pencil and active contacts}

Define the reversal matrix
\[
\Pi_n e_j=e_{n+1-j}.
\]
For real \(x\), define the symmetric reversal pencil
\[
\mathbf G_n(x,a)
:=
(x\mathbf I_n-\mathbf A_n(a))\Pi_n.
\]
Since \(\Pi_n\) is orthogonal,
\[
s_{\min}(x\mathbf I_n-\mathbf A_n(a))
=
\operatorname{dist}(0,\sigma(\mathbf G_n(x,a))).
\]

\begin{definition}[Simple interior active contact]
A triple
\[
(x,\lambda,w)
\]
is a simple interior active contact for length \(n\) if
\[
\mathbf G_n(x,a)w=\lambda w,\qquad \|w\|=1,
\]
\[
|\lambda|=\eta_n(a),
\]
the eigenvalue \(\lambda\) is simple and is the unique eigenvalue of \(\mathbf G_n(x,a)\)
closest to \(0\), and \(x\) lies in the interior of the real spectral hull.
\end{definition}

Let
\[
E_n=\operatorname{diag}(0,1,\dots,n-1),
\]
\[
M_n(w)=w^TE_nw,
\]
and
\[
Q_n(w)=2M_n(w)-(n-1).
\]

\begin{lemma}[Krein neutrality]
\label{lem:neutrality}
At a simple interior active contact,
\[
w^T\Pi_nw=0.
\]
\end{lemma}

\begin{proof}
Because
\[
\partial_x\mathbf G_n(x,a)=\Pi_n,
\]
Hellmann's formula gives
\[
\partial_x\lambda=w^T\Pi_nw.
\]
At an interior maximum of \(|\lambda(x)|\),
\[
\partial_x|\lambda|=0.
\]
Since \(\lambda\ne0\), this gives \(w^T\Pi_nw=0\).
\end{proof}

\begin{lemma}[Orientation derivative]
\label{lem:orientation-derivative}
Let
\[
\tau=-\frac12\log a.
\]
Along a differentiable simple active branch,
\[
\partial_\tau\log|\lambda|
=
-1-Q_n(w).
\]
Equivalently,
\[
\partial_\tau\log|\lambda|
=
n-2-2M_n(w).
\]
\end{lemma}

\begin{proof}
Let
\[
D_n(\tau)=\operatorname{diag}(1,e^{-\tau},\dots,e^{-(n-1)\tau}).
\]
Using the diagonal symmetrization of \(\mathbf A_n(a)\), one computes
\[
\partial_\tau \mathbf G_n
=
(n-2)\mathbf G_n-E_n\mathbf G_n-\mathbf G_nE_n.
\]
Hellmann's formula gives
\[
\partial_\tau\lambda
=
w^T(\partial_\tau\mathbf G_n)w.
\]
Since \(\mathbf G_nw=\lambda w\),
\[
\partial_\tau\lambda
=
(n-2)\lambda-2\lambda w^TE_nw.
\]
Dividing by \(\lambda\) gives the result.
\end{proof}

At a smooth crossing
\[
\eta_{n+1}(a)=\eta_n(a),
\]
with simple active contacts \(w_n,w_{n+1}\), one has
\[
\partial_\tau\log\frac{\eta_{n+1}}{\eta_n}
=
Q_n(w_n)-Q_{n+1}(w_{n+1}).
\]
Thus an upward crossing requires
\[
Q_{n+1}(w_{n+1})\le Q_n(w_n).
\]

\section{Tail-cone constraint}

For \(1\le k\le\lfloor n/2\rfloor\), define
\[
\Delta_{n,k}(w)
=
\sum_{j=1}^{k}\left(w_{n+1-j}^2-w_j^2\right).
\]

\begin{theorem}[Active tail-cone membership]
\label{thm:tail-cone}
At every simple interior active contact,
\[
\Delta_{n,k}(w)\ge0,
\qquad
k=1,\dots,\left\lfloor\frac n2\right\rfloor.
\]
\end{theorem}

\begin{proof}[Proof sketch]
Approach the real active contact from the upper half-plane. Let
\[
z=x+iy,\qquad y>0,
\]
and let \(v(y)\) be the right singular vector in the canonical gauge
\[
(z\mathbf I-\mathbf A)v=s\,u,\qquad u=\Pi_n\overline v.
\]
Define the current
\[
W_j(y)
=
v_{j+1}(y)v_{n+1-j}(y)-a\,v_j(y)v_{n-j}(y),
\qquad j=0,\dots,n,
\]
with \(v_0=v_{n+1}=0\). The singular-vector current identity gives
\[
W_j-W_{j-1}
=
s\left(|v_j|^2-|v_{n+1-j}|^2\right).
\]
The phase-continuation argument in the upper half-plane gives \(W_j>0\) for
\(1\le j\le n-1\). Therefore
\[
\sum_{j=1}^{k}|v_j|^2
>
\sum_{j=1}^{k}|v_{n+1-j}|^2.
\]
Letting \(y\downarrow0\) and using \(w=\Pi_nv(0)\), up to sign, gives the result.
\end{proof}

\begin{corollary}
\[
Q_n(w)\ge0
\]
at every simple interior active contact.
\end{corollary}

\begin{proof}
By Abel summation,
\[
Q_n(w)
=
\sum_{j=1}^{\lfloor n/2\rfloor}(n+1-2j)
\left(w_{n+1-j}^2-w_j^2\right),
\]
and the right side is nonnegative by Theorem~\ref{thm:tail-cone}.
\end{proof}

\section{Paired Weyl formulation}

We record the even-to-odd Weyl formulation. The odd-to-even version remains a
separate unfinished part of the proof.

Let \(n=2m\). Define paired variables
\[
y_k=
\binom{\ell_k}{r_k}
:=
\binom{w_k}{w_{2m+1-k}},
\qquad k=1,\dots,m.
\]
Then the reversal eigenproblem is equivalent to
\[
C y_{k+1}+xJ y_k+C^T y_{k-1}=\lambda y_k,
\qquad k=1,\dots,m,
\]
with
\[
J=
\begin{pmatrix}0&1\\1&0\end{pmatrix},
\qquad
C=
\begin{pmatrix}0&-1\\-a&0\end{pmatrix}.
\]
The left boundary is
\[
y_0=0.
\]
The even terminal condition is
\[
y_{m+1}=Jy_m.
\]
Equivalently,
\[
C y_{m+1}=B_Ey_m,
\qquad
B_E=
\begin{pmatrix}-1&0\\0&-a\end{pmatrix}.
\]

For odd length \(2m+1\), one has paired variables
\[
\widetilde y_k=
\binom{\widetilde\ell_k}{\widetilde r_k},
\qquad k=1,\dots,m,
\]
and a center coordinate \(c\). The center equation is
\[
(\widetilde x-\widetilde\lambda)c
=
\widetilde\ell_m+a\widetilde r_m.
\]
Eliminating \(c\) gives the odd terminal boundary matrix
\[
B_O(\widetilde x,\widetilde\lambda)
=
-\frac{1}{\widetilde x-\widetilde\lambda}
\binom{1}{a}\binom{1}{a}^T,
\]
with the non-negotiable condition
\[
\widetilde x-\widetilde\lambda\ne0.
\]

\section{Mixed Cauchy decomposition}

Let \(Y=(y_1,\dots,y_m)\) be the paired even vector and
\[
\widetilde W=(\widetilde Y,c),
\qquad
\widetilde Y=(\widetilde y_1,\dots,\widetilde y_m),
\]
be the paired odd vector plus its center.

Let
\[
\mathcal J=\operatorname{diag}(J,\dots,J).
\]
Even stationarity is
\[
\langle \mathcal JY,Y\rangle=0.
\]
Odd stationarity is
\[
\langle \mathcal J\widetilde Y,\widetilde Y\rangle=-c^2.
\]

Let
\[
P_\pm=\frac{I\pm\mathcal J}{2}.
\]
Write
\[
Y_\pm=P_\pm Y,
\qquad
\widetilde Y_\pm=P_\pm\widetilde Y.
\]
Then
\[
\|Y_+\|^2=\|Y_-\|^2=\frac12,
\]
and
\[
\|\widetilde Y_+\|^2=\frac{1-2c^2}{2},
\qquad
\|\widetilde Y_-\|^2=\frac12.
\]

Define
\[
\mathcal N=\langle \widetilde Y,Y\rangle,
\qquad
\mathcal S=\langle \widetilde Y,\mathcal JY\rangle,
\]
and
\[
\alpha=\mathcal N+\mathcal S,
\qquad
\beta=\mathcal N-\mathcal S.
\]
Decompose
\[
\widetilde Y_+=\alpha Y_+ + U_+,
\qquad
U_+\perp Y_+,
\]
\[
\widetilde Y_-=\beta Y_- + U_-,
\qquad
U_-\perp Y_-.
\]
The Cauchy defects are
\[
\mathcal D_+=2\|U_+\|^2,
\qquad
\mathcal D_-=2\|U_-\|^2.
\]
The exact-rigid branch is
\[
\mathcal D_+=\mathcal D_-=0.
\]
The non-rigid branch is
\[
\mathcal D_++\mathcal D_->0.
\]

\section{Exact-rigid even-to-odd branch}

This section records the current proof that the exact-rigid even-to-odd branch is
empty.

\subsection{Exact-rigid form}

Exact rigidity means
\[
\widetilde Y=\alpha Y_+ +\beta Y_-.
\]
The alignment-sign argument gives, after changing the global sign of the odd vector,
\[
\alpha\beta>0.
\]
Hence one may write
\[
\widetilde y_k=L_\rho y_k,
\qquad
L_\rho=\rho P_+ +P_-,
\qquad
0\le\rho\le1,
\]
where now \(P_\pm=(I\pm J)/2\) on each block, and
\[
\rho^2=1-2c^2.
\]
The case \(c=0\), equivalently \(\rho=1\), is impossible for an upward crossing. Thus
the nontrivial exact-rigid case has
\[
0<\rho<1.
\]

\subsection{Flat-tail exclusion}

\begin{lemma}[Flat-tail exact rigidity is impossible]
If
\[
\Delta_m=\sum_{k=1}^{m}(r_k^2-\ell_k^2)=0,
\]
then no exact-rigid active crossing exists.
\end{lemma}

\begin{proof}
The even terminal current identity gives
\[
\lambda\Delta_m=\ell_m^2-ar_m^2.
\]
Thus
\[
\ell_m^2=ar_m^2.
\]
Since the terminal block cannot vanish, \(r_m\ne0\), so
\[
\ell_m/r_m=\pm\sqrt a.
\]
Compatibility with the odd center and current equations forces
\[
\ell_m/r_m=\sqrt a,
\qquad
\rho=\frac{1+a}{(1+\sqrt a)^2}.
\]
But then
\[
r_m^2-\ell_m^2=(1-a)r_m^2>0.
\]
Since
\[
\Delta_m=\Delta_{m-1}+(r_m^2-\ell_m^2)=0,
\]
we obtain
\[
\Delta_{m-1}<0,
\]
contradicting the tail-cone condition.
\end{proof}

\subsection{\texorpdfstring{Same-sign exact rigidity for \(m\ge4\)}{Same-sign exact rigidity for m >= 4}}

Assume
\[
\widetilde\lambda=\lambda.
\]
The local current law becomes the projective relation
\[
(\ell_k-r_k)(\ell_{k+1}-r_{k+1})
+
\rho(\ell_k+r_k)(\ell_{k+1}+r_{k+1})=0.
\]
In Cayley coordinate
\[
z_k=\frac{\ell_k-r_k}{\ell_k+r_k},
\]
this is
\[
z_kz_{k+1}=-\rho.
\]
Thus the projective lines satisfy
\[
[y_{k+2}]=[y_k].
\]

\begin{theorem}[Same-sign exact-rigid exclusion for \(m\ge4\)]
\label{thm:same-sign-mge4}
There is no same-sign exact-rigid even-to-odd active crossing for \(m\ge4\).
\end{theorem}

\begin{proof}
Since \([y_{k+2}]=[y_k]\), write
\[
y_3=b\,y_1,\qquad y_4=g\,y_2.
\]
The recurrence at \(k=1\) is
\[
Cy_2+(xJ-\lambda I)y_1=0.
\]
The recurrence at \(k=3\) is
\[
Cy_4+(xJ-\lambda I)y_3+C^Ty_2=0.
\]
Substituting \(y_3=by_1\), \(y_4=gy_2\), and using the \(k=1\) recurrence gives
\[
((g-b)C+C^T)y_2=0.
\]
Since
\[
(g-b)C+C^T
=
\begin{pmatrix}
0&-(g-b)-a\\
-a(g-b)-1&0
\end{pmatrix},
\]
and \(0<a<1\), this forces
\[
y_2\parallel e_1
\quad\text{or}\quad
y_2\parallel e_2.
\]
If \(y_2\parallel e_2\), the projective law implies
\[
\frac{\ell_1}{r_1}=\frac{1+\rho}{1-\rho}>1,
\]
contradicting \(\Delta_1\ge0\). If \(y_2\parallel e_1\), the recurrence at \(k=1\) and
\(k=2\) gives two scalar equations whose signs are incompatible for \(0<a<1\).
Thus both coordinate cases are impossible.
\end{proof}

\subsection{\texorpdfstring{Opposite-sign exact rigidity for \(m\ge4\)}{Opposite-sign exact rigidity for m >= 4}}

Assume
\[
\widetilde\lambda=-\lambda.
\]
In \(J\)-eigen-coordinates, the exact-rigid modulation equations imply
\[
y_{k+1}=y_{k-1}+Ry_k
\]
for a fixed \(2\times2\) matrix \(R\). Combining this with the even recurrence gives
\[
y_{k-1}=Ty_k
\]
for a fixed \(2\times2\) matrix \(T\).

\begin{theorem}[Opposite-sign exact-rigid exclusion for \(m\ge4\)]
\label{thm:opposite-mge4}
There is no opposite-sign exact-rigid even-to-odd active crossing for \(m\ge4\).
\end{theorem}

\begin{proof}
Because \(y_0=0\),
\[
Ty_1=0.
\]
For \(m\ge4\), the identities at \(k=2,3\) give
\[
TSy_1=y_1,
\qquad
TSy_2=y_2,
\]
where \(S=T+R\). If \(y_1,y_2\) are linearly independent, then \(TS=I\), contradicting
\(\det T=0\). Hence \(y_2=\kappa y_1\). Then all blocks satisfy
\[
y_k=\kappa^{k-1}y_1.
\]
Even stationarity gives
\[
0=\sum_{k=1}^{m}y_k^TJy_k
=
\left(\sum_{k=1}^{m}\kappa^{2(k-1)}\right)y_1^TJy_1.
\]
Hence \(y_1^TJy_1=0\), so \(y_1\parallel e_1\) or \(y_1\parallel e_2\) in the original
coordinate axes. The \(k=1\) recurrence then forces \(\lambda=0\), impossible at an
active barrier contact.
\end{proof}

\subsection{Low exact-rigid cases}

The low cases \(m=1,2,3\) were handled separately.

\begin{proposition}[\(m=1\)]
There is no exact-rigid active crossing for \(m=1\).
\end{proposition}

\begin{proof}
Here \(y_1=(\ell,r)^T\). Even stationarity gives \(\ell r=0\). The branch \(y_1\parallel e_1\)
violates the tail cone. The branch \(y_1\parallel e_2\) gives
\[
x=0,\qquad \lambda=-a.
\]
Exact rigidity gives
\[
\widetilde y_1=L_\rho e_2
=
\binom{(\rho-1)/2}{(1+\rho)/2}.
\]
Substitution into the odd terminal equation, with
\[
\widetilde x-\widetilde\lambda\ne0,
\]
gives the necessary polynomial equation
\[
P_1(a,\rho)=0,
\]
where
\[
\begin{aligned}
P_1(a,\rho)
={}&
-a^3\rho^3-3a^3\rho^2+5a^3\rho-a^3
+a^2\rho^3+a^2\rho^2+7a^2\rho-a^2\\
&+a\rho^3-a\rho^2-a\rho+a-\rho^3+3\rho^2-3\rho+1.
\end{aligned}
\]
Its Bernstein coefficients in \(\rho\in[0,1]\) are
\[
(1-a)(1+a)^2,\qquad
\frac{2a(1+a)^2}{3},\qquad
\frac{4a^2(a+3)}{3},\qquad
8a^2,
\]
all positive for \(0<a<1\). Thus \(P_1>0\), contradiction.
\end{proof}

\begin{proposition}[\(m=2\)]
There is no exact-rigid active crossing for \(m=2\).
\end{proposition}

\begin{proof}[Proof summary]
This is the transition \(4\to5\). In the finite chart write
\[
y_1=\binom{s}{1},
\qquad
y_2=h\binom{t}{1}.
\]
Even stationarity gives
\[
t=-\frac{s}{h^2}.
\]
The even recurrence and terminal condition imply
\[
q:=\frac hs
\]
satisfies
\[
aq^2+2(1+a)q-1=0.
\]
The positive branch \(q=q_+(a)\in(0,1/2)\) contradicts \(\Delta_2\ge0\). Thus
\[
q=q_-(a)<-2-\sqrt5.
\]
Write \(q=-p\), so \(p>2+\sqrt5\). Then
\[
a=\frac{1+2p}{p(p-2)}.
\]
The same-sign branch gives
\[
\rho=
\frac{(1-s)(1+p^2s)}{(1+s)(1-p^2s)}.
\]
The even equations imply
\[
|s|>\frac1{p^2}.
\]
If \(s>0\), then \(1-p^2s<0\), so \(\rho<0\). If \(s<0\), then
\(1+p^2s<0\), again \(\rho<0\). Hence the same-sign branch is impossible.

For the opposite-sign finite chart, denominator-preserving elimination gives the
necessary condition
\[
\mathcal P_{42}(q)=0,
\qquad q<-2-\sqrt5,
\]
where \(\mathcal P_{42}\) is the exact residual polynomial obtained from the even
constraints, the odd bulk equation, and one denominator-preserving component of
the odd terminal equation. Under the substitution
\[
q=-(2+\sqrt5+u),\qquad u>0,
\]
one obtains
\[
\mathcal P_{42}(-(2+\sqrt5+u))
=
\sum_{j=0}^{42}(A_j+B_j\sqrt5)u^j
\]
with
\[
A_j>0,\qquad B_j>0\qquad(0\le j\le42).
\]
Thus \(\mathcal P_{42}>0\) on the admissible interval, contradiction.

The projective singular charts \(y_2\parallel e_1\) and \(y_2\parallel e_2\) are excluded
directly. In the \(e_2\)-chart one obtains \(\lambda=0\). In the \(e_1\)-chart the
remaining terminal residual reduces to a strictly positive expression after setting
\(a=t^2\in(0,1)\). Hence no branch remains.
\end{proof}

\begin{proposition}[\(m=3\), same sign]
There is no same-sign exact-rigid active crossing for \(m=3\).
\end{proposition}

\begin{proof}
Use \(J\)-eigen-coordinates
\[
y_k=u_ke_+ +v_ke_-.
\]
The exact-rigid same-sign modulation equations give, for some \(\tau\),
\[
u_2=\tau v_1,\qquad
v_2=-\rho\tau u_1,
\]
and
\[
u_3=(1-\rho\tau^2)u_1,\qquad
v_3=(1-\rho\tau^2)v_1.
\]
Thus
\[
y_3=(1-\rho\tau^2)y_1.
\]
The cases \(u_1=0\) or \(v_1=0\) force \(\tau=0\), then \(y_2=0\), and then the
recurrence implies \(y_1=0\), impossible. Thus \(u_1v_1\ne0\). Setting
\(q=v_1/u_1\), the \(k=1\) recurrence eliminates \(x,\lambda\); substituting into the
\(k=2\) recurrence gives
\[
\tau^2(q^2+\rho)=2,
\qquad
q^2=\rho.
\]
Therefore \(\rho\tau^2=1\), hence \(y_3=0\). Since the terminal condition gives
\(y_4=Jy_3=0\), two consecutive zero blocks force the trivial solution, contradiction.
\end{proof}

\begin{proposition}[\(m=3\), opposite sign]
There is no opposite-sign exact-rigid active crossing for \(m=3\).
\end{proposition}

\begin{proof}
Use \(J\)-eigen-coordinates, in which
\[
C=
\begin{pmatrix}
-\mu&\nu\\
-\nu&\mu
\end{pmatrix},
\qquad
\mu=\frac{1+a}{2},\quad
\nu=\frac{1-a}{2}.
\]
Thus \(\mu>\nu>0\). The opposite-sign modulation gives parameters \(A,B\) such that
\[
u_2=-Bv_1,\qquad
v_2=-\rho A u_1,
\]
and
\[
u_3=(1+\rho AB)u_1,\qquad
v_3=(1+\rho AB)v_1.
\]
If \(u_1=0\), then the \(k=1\) recurrence forces \(B=0\), and the \(k=2\) recurrence
forces \(2\mu=0\), impossible. Hence \(u_1\ne0\). Normalize \(u_1=1\) and set
\(q=v_1\). The \(k=1\), \(k=2\), and \(k=3\) recurrences reduce to
\[
A\bigl(2\mu\rho+\nu(1+\rho)q\bigr)
=
Bq\bigl(2\mu q+\nu(1+\rho)\bigr),
\]
\[
B^2q^2-AB\rho-2=0,
\]
\[
\nu\rho(1+\rho)A(A-B)+4\mu q=0,
\]
and finally
\[
\mu q(5\rho+1)+\nu\rho(1+\rho)=0,
\]
\[
q\bigl(4\mu^2\rho+\nu^2(1+\rho)\bigr)+\mu\nu\rho(1+\rho)=0.
\]
The first of the last two equations gives
\[
q=-\frac{\nu\rho(1+\rho)}{\mu(5\rho+1)}.
\]
Substituting this into the second gives
\[
\frac{\nu\rho(1+\rho)^2(\mu^2-\nu^2)}
{\mu(5\rho+1)}=0.
\]
But
\[
\mu^2-\nu^2=a>0,
\]
and all other factors are positive. Contradiction.
\end{proof}

\begin{theorem}[Exact-rigid even-to-odd branch eliminated]
\label{thm:rigid-eliminated}
There is no exact-rigid simple interior active even-to-odd upward crossing.
\end{theorem}

\begin{proof}
The flat-tail branch is impossible. The low block lengths \(m=1,2,3\) are impossible.
For \(m\ge4\), both the same-sign and opposite-sign branches are impossible by
Theorems~\ref{thm:same-sign-mge4} and \ref{thm:opposite-mge4}.
\end{proof}

\section{The non-rigid even-to-odd branch}

The exact-rigid branch is eliminated, so the remaining simple even-to-odd branch
has
\[
\mathcal D_++\mathcal D_->0.
\]

Define
\[
Z=(Z_Y,c),
\qquad
Z_Y=\alpha Y_++\beta Y_-,
\]
and
\[
U=(U_++U_-,0).
\]
Then
\[
\widetilde W=Z+U.
\]
Let
\[
d=\|U\|^2
=
\frac{\mathcal D_++\mathcal D_-}{2}.
\]
Then
\[
d>0
\]
in the non-rigid branch.

Let
\[
\widetilde G=
\mathbf G_{2m+1}(\widetilde x,a),
\qquad
\widetilde G\widetilde W=\widetilde\lambda\widetilde W,
\qquad
\|\widetilde W\|=1.
\]
Define the squared-Rayleigh excess
\[
\mathcal Q(v)
=
\|\widetilde Gv\|^2-\widetilde\lambda^2\|v\|^2.
\]
Because \(|\widetilde\lambda|\) is the least absolute eigenvalue of \(\widetilde G\),
\[
\mathcal Q(v)\ge0.
\]

\begin{lemma}[Squared-Rayleigh identity]
\label{lem:QZQU}
One has
\[
\mathcal Q(Z)=\mathcal Q(U)\ge0.
\]
Moreover, if the active odd eigenvalue is simple and
\[
g_{\min}
=
\min_{\nu\in\sigma(\widetilde G),\,\nu\ne\widetilde\lambda}
(\nu^2-\widetilde\lambda^2)>0,
\]
then
\[
\mathcal Q(Z)\ge g_{\min}d(1-d).
\]
\end{lemma}

\begin{proof}
Since
\[
(\widetilde G^2-\widetilde\lambda^2I)\widetilde W=0
\]
and \(\widetilde W=Z+U\), the bilinear form associated with \(\mathcal Q\) gives
\[
\mathcal Q(Z)=\mathcal Q(U).
\]
Furthermore,
\[
\langle Z,\widetilde W\rangle=1-d,
\qquad
\|Z\|^2=1-d.
\]
Thus, writing
\[
Z=(1-d)\widetilde W+Z_\perp,\qquad Z_\perp\perp\widetilde W,
\]
we get
\[
\|Z_\perp\|^2=d(1-d).
\]
The spectral theorem then gives
\[
\mathcal Q(Z)\ge g_{\min}d(1-d).
\]
\end{proof}

\begin{lemma}[Stationarity defect of the rigid projection]
\label{lem:stationarity-defect}
The odd stationarity defect of \(Z\) is
\[
\langle \mathcal JZ_Y,Z_Y\rangle+c^2
=
\frac{\mathcal D_--\mathcal D_+}{2}.
\]
\end{lemma}

\begin{proof}
Since
\[
Z_Y=\alpha Y_++\beta Y_-,
\]
we have
\[
\langle \mathcal JZ_Y,Z_Y\rangle
=
\frac{\alpha^2-\beta^2}{2}.
\]
The defect identities are
\[
\alpha^2+\mathcal D_+=1-2c^2,
\qquad
\beta^2+\mathcal D_-=1.
\]
Thus
\[
\alpha^2-\beta^2=\mathcal D_--\mathcal D_+-2c^2,
\]
and the formula follows.
\end{proof}

The exact orientation identity is
\[
\begin{aligned}
Q_{2m}(Y)-Q_{2m+1}(\widetilde Y)
&=
(1-\alpha\beta)Q_{2m}(Y)
-\alpha\beta\,\Delta_m(Y)
\\
&\quad
-2\alpha\langle Y_+,W_OHU_-\rangle
\\
&\quad
-2\beta\langle U_+,W_OHY_-\rangle
\\
&\quad
-2\langle U_+,W_OHU_-\rangle,
\end{aligned}
\]
where \(H=\operatorname{diag}(-1,1)\) on each pair, and \(W_O\) is the diagonal orientation
weight for the odd length.

An upward crossing requires
\[
Q_{2m+1}(\widetilde Y)\le Q_{2m}(Y),
\]
or
\[
Q_{2m}(Y)-Q_{2m+1}(\widetilde Y)\ge0.
\]

The non-rigid branch is not yet excluded. The missing theorem is a defect-vector
Picone identity proving that the last three defect terms force
\[
Q_{2m+1}(\widetilde Y)>Q_{2m}(Y)
\]
whenever
\[
\mathcal D_++\mathcal D_->0.
\]

\section{Current status}

The connectedness threshold theorem is not fully proved yet. The current proof
establishes the following.

\begin{enumerate}[label=\textup{(\arabic*)}]
\item The uploaded manuscript proves the barrier criterion
\[
\Sigma_\varepsilon^{(n)}\text{ connected}
\iff
\varepsilon>\eta_n(a).
\]

\item Therefore the threshold problem is reduced to
\[
\eta_{n+1}(a)\le\eta_n(a).
\]

\item The uniform small-\(a\) strip
\[
0<a\le10^{-2}
\]
is excluded for all \(n\ge2\).

\item The perturbative near-Hermitian layer
\[
1-a\le \frac{1}{4(n+2)^2}
\]
is excluded.

\item Large-\(n\) compact-middle and mesoscopic near-Hermitian asymptotics reduce
possible counterexamples to a finite middle block.

\item In the finite middle block, the simple even-to-odd active-crossing problem has
been reduced to rigid and non-rigid Weyl-contact branches.

\item The exact-rigid simple even-to-odd branch is eliminated completely.
\end{enumerate}

\section{Unfinished gaps}

\begin{gap}[Non-rigid simple even-to-odd branch]
The main remaining gap inside the current Weyl-contact route is the non-rigid branch:
\[
\mathcal D_++\mathcal D_->0.
\]
The missing result is a defect-vector Picone identity showing that non-rigidity forces
\[
Q_{2m+1}(\widetilde Y)>Q_{2m}(Y).
\]
This would contradict the upward-crossing condition
\[
Q_{2m+1}(\widetilde Y)\le Q_{2m}(Y).
\]
At present we only have the squared-Rayleigh lower bound
\[
\mathcal Q(Z)\ge g_{\min}d(1-d),
\]
and the stationarity-defect identity
\[
\langle \mathcal JZ_Y,Z_Y\rangle+c^2
=
\frac{\mathcal D_--\mathcal D_+}{2}.
\]
These do not yet close the non-rigid branch.
\end{gap}

\begin{gap}[Odd-to-even analogue]
The Weyl-contact analysis above is written for even-to-odd transitions
\[
2m\to2m+1.
\]
The odd-to-even transition
\[
2m+1\to2m+2
\]
still needs its own paired Weyl formulation, exact-rigid elimination, and non-rigid
defect analysis.
\end{gap}

\begin{gap}[Degenerate active contacts]
The current Weyl proof assumes simple interior active contacts. One still has to pass
to nonsimple active eigenvalues, endpoint contacts, and nonsmooth maxima. This should
be possible by semialgebraic stratification, limiting arguments, or perturbation of
simple active branches, but it is not yet written.
\end{gap}

\begin{gap}[Global crossing closure]
Even after excluding simple upward crossings, one must write the global crossing
argument carefully. Namely, if
\[
R_n(a)=\frac{\eta_{n+1}(a)}{\eta_n(a)}
\]
ever exceeds \(1\), then there must be an upward crossing in a suitable stratum.
This needs a formal stratified argument over the finite middle block.
\end{gap}

\begin{gap}[Full rigor of asymptotic reductions]
The large-\(n\) compact-middle and mesoscopic transition estimates are based on the
exact Green kernel and Volterra scaling. A final manuscript must include all uniform
operator-norm and Euler--Maclaurin remainder estimates, especially for
\[
1\le \Gamma\le C\log n.
\]
\end{gap}

\begin{gap}[Finite residual audit]
The \(m=2\) exact-rigid proof uses an exact degree-\(42\) residual polynomial with a
positive-coefficient expansion after the substitution
\[
q=-(2+\sqrt5+u).
\]
This is an exact algebraic proof, but it should be audited and placed in a separate
appendix with all coefficients or a shorter hand factorization if one can be found.
\end{gap}

\end{document}
